\section{A Cartan bound for fixed root data}\label{sec:geometry}

In this section the characteristic of the algebraic group is denoted by
\(r\), and the coefficient characteristic is the fixed prime \(p\ne r\).
A root datum includes its character and cocharacter lattices, and hence
the central torus as well as the semisimple roots.

\begin{theorem}[Fixed-root Cartan bound]\label{thm:geometric-cartan}
Fix a finite collection \(\mathcal R\) of root data and a finite collection
of the usual pinned diagram automorphisms and exceptional diagram
isogenies on these data. There is a constant \(C=C(\mathcal R,p,M)\), also
depending on this fixed collection of maps, with the following property.
Let \(\mathbf G\) be a connected reductive group over
\(\overline{\F}_r\), with root datum in \(\mathcal R\), and let \(F\) be a
Steinberg endomorphism which, after an isomorphism of rational structures,
has the form
\[
 F=\theta\operatorname{Fr}_Q.
\]
Here \(Q\) is a power of \(r\), \(\operatorname{Fr}_Q\) is split
Frobenius, and \(\theta\) belongs to the specified finite collection;
exceptional isogenies are allowed only in their prescribed defining
characteristics. Put \(\Gamma=\mathbf G^F\).
For every subgroup \(Z\le Z(\Gamma)\) and every block \(b\) of
\(k[\Gamma/Z]\) with defect groups of order at most \(M\), all entries of
the Cartan matrix of \(b\) are at most \(C\).
\end{theorem}

The assertion places no bound on a Sylow \(p\)-subgroup of \(\Gamma\)
or on the \(p\)-part of a rational maximal torus. We will construct
projective modules using the horocycle correspondence, and bound their
characters only after projection to the prescribed quotient block.
The horocycle category and the passage from tilting sheaves to
projective representations follow Eteve's construction
\cite[Lemma~2.1.1 and Theorems~2.3.2, 3.1.1(ii), 3.2.4]{Eteve}.
Degree-zero projectivity of the corresponding tilting images was
also obtained independently by Zhu
\cite[Theorem~4.91]{ZhuTame}.  The uniform estimates needed here
are proved below: they must be independent of the rational torus
and its monodromy parameter.

We use the following standard parts of Deligne--Lusztig theory for
connected reductive groups with a Steinberg endomorphism, including its
isogeny form, in the large-parameter range used below: torus duality,
the partition into geometric Lusztig series,
the character formula, orthogonality, and the uniform Steinberg formula.
Their relevant forms are recalled when used; see
\cite[Theorems~4.2, 6.2 and 6.8, Proposition~5.22,
Corollaries~6.3 and 7.14, and Section~11]{DL}.
All constructible sheaves have coefficient characteristic different from
the geometric characteristic. Calculations with \(k\)-coefficients may
be descended to a sufficiently large finite subfield of \(k\), and
integral calculations to a finite extension of a complete splitting
discrete valuation ring. Tate twists do not affect any dimension below.

We use constructible descent and the six operations with finite
coefficients, and their adic counterparts, as in
\cite{LaszloOlssonFinite,LaszloOlssonAdic}.  For the finite local
coefficient algebras below, the assertions needed in
Lemma~\ref{geom:flat-costalk} are obtained after forgetting to a
finite coefficient field, as in its proof; no self-duality assertion
for an arbitrary finite local algebra is used.

\subsection{Strata and perfect extraordinary restrictions}

The first task is qualitative: the restrictions needed to glue
standard objects must be perfect complexes over their stabilizer
algebras.  Once this is established, a separate estimate will bound
the number of their indecomposable projective terms independently
of the stabilizer orders.

Work with one root datum. Choose an \(F\)-stable Borel pair
\(\mathbf T\subset\mathbf B=\mathbf T\mathbf U\), write \(W\) for the
Weyl group, and let \(\ell\) be its length function. A bounded range of
\(Q\) gives only finitely many groups for the specified data. It may
therefore be omitted until the end. In particular, assume \(dF=0\).
For \(x\in W\), choose \(\dot x\in N_{\mathbf G}(\mathbf T)\), and put
\[
 \mathbf G_x=\mathbf B\dot x\mathbf B,
 \qquad C_x=\mathbf G_x/\mathbf B,
 \qquad A_x=\mathbf T^{xF}.
\]
The torus Lang map is an \'etale isogeny, so \(A_x\) is a finite abelian
group of order prime to \(r\). It identifies with the rational points
of a corresponding \(F\)-stable maximal torus \(T_x\) of \(\mathbf G\).

Define
\[
 \mathcal H=
 \{(ut,u'F(t)):u,u'\in\mathbf U,\ t\in\mathbf T\},
 \qquad \mathcal X=[\mathbf G/\mathcal H],
\]
where \((b,b')\) sends \(g\) to \(bg(b')^{-1}\).
Its strata are
\(j_x:\mathcal X_x=[\mathbf G_x/\mathcal H]\hookrightarrow\mathcal X\).
The torus Lang isogeny makes this action transitive on each stratum.
The stabilizer of \(\dot x\) is
\[
 (\mathbf U\cap{}^{\dot x}\mathbf U)\rtimes A_x.
\]
Indeed its torus equation is \(t=xF(t)\), and its unipotent equation is
\(u={} ^{\dot x}u'\). The section of the torus quotient is
\(t\mapsto(t,F(t))\).

Consequently ordinary inflation, without a shift, identifies
\begin{equation}\label{geom:stratum-category}
 D^b_c(\mathcal X_x,k)\simeq D^b(kA_x\text{-mod}).
\end{equation}
Here and throughout, constructibility includes finite-dimensional
cohomology. To verify the identification, cohomology sheaves on the
classifying stack are representations of the component group \(A_x\).
The components in the nerve of the stabilizer are affine spaces, since
its connected unipotent radical is split. Their positive ordinary
\'etale cohomology vanishes. Descent therefore computes Ext by the
ordinary group-cohomology complex of \(A_x\). Truncation proves the
bounded derived assertion. This also explains why the equivalence
retains modular extensions even when \(p\mid |A_x|\); compare
\cite[Lemma 2.1.1]{Eteve}.

For a character \(\phi:A_x\to k^\times\), let \(L_{x,\phi}\) be its
one-dimensional module and let \(E_{x,\phi}\) be its projective cover.
Writing \(A_x=(A_x)_p\times(A_x)_{p'}\), we have
\[
 E_{x,\phi}=k[(A_x)_p]\otimes_k L_{x,\phi}.
\]
Set
\[
 \Delta_{x,\phi}=j_{x,!}E_{x,\phi}[\ell(x)],
 \qquad
 \nabla_{x,\phi}=Rj_{x,*}E_{x,\phi}[\ell(x)].
\]
A standard, respectively costandard, flag is a finite filtration by
extensions with factors of the displayed standard, respectively
costandard, objects, without additional shifts.

Pullback to \(\mathbf G\), followed by \([\dim\mathbf B]\), makes both
objects perverse. The cell has dimension
\(\dim\mathbf B+\ell(x)\), its coefficient is lisse, and its inclusion
is affine and quasi-finite: both \(\mathbf G_x\) and \(\mathbf G\) are
affine. Apply affine quasi-finite perverse exactness
\cite[Corollary 4.1.3]{BBD}. These assertions hold on invariant open
unions of cells as well, by base change.

We first prove the coefficient lemma that supplies perfection.

\begin{lemma}\label{geom:flat-costalk}
Let \(R\) be a finite-dimensional commutative local algebra over a
finite field of characteristic \(p\), with residue field \(k_0\).
Let \(Y\) be a variety over an algebraically closed field of
characteristic different from \(p\), and let \(\mathcal D\) be a
constructible sheaf of \(R\)-modules with projective geometric stalks.
For a locally closed immersion \(i:Z\hookrightarrow Y\) and a geometric
point \(z\in Z\), the complex \((Ri^!\mathcal D)_z\) is perfect over
\(R\). The same statement holds after extension to algebraic coefficient
fields when the data descend to finite coefficients.
\end{lemma}

\begin{proof}
Put \(P=(Ri^!\mathcal D)_z\). Forgetting the coefficient action commutes
with extraordinary restriction: the forgetful functor preserves
injectives because its left adjoint is exact. Constructible finiteness
therefore gives bounded finite \(R\)-cohomology for \(P\).
Projectivity on stalks says that \(\mathcal D\) is flat over \(R\).

For any bounded complex \(B\) of finite free \(R\)-modules, adjunction
with extension by zero gives the projection comparison
\[
 P\otimes_R B\xrightarrow{\sim}
 \bigl(Ri^!(\mathcal D\otimes_R B)\bigr)_z.
\]
It is an isomorphism for \(B=R\), hence for finite sums, shifts and
cones. To extend this comparison to the residue field, take a
bounded-above finite free resolution of \(k_0\), and let \(B_n\) be
its brutal truncation in degrees \([-n,0]\). There is a finite module
\(N_n\) such that
\(\operatorname{Cone}(B_n\to k_0)\simeq N_n[n+1]\).
Choose \(b\) with \(P\in D^{\le b}(R)\). Right \(t\)-exactness of
derived tensor gives
\[
 P\otimes_R^L N_n[n+1]\in D^{\le b-n-1}(R).
\]
Choose also an upper cohomological bound \(d\) for
\((Ri^!(\mathcal D\otimes_R k_0))_z\).
Every finite \(R\)-module has a finite filtration by \(k_0\).
Flatness and the long exact cohomology sequence show that
\((Ri^!(\mathcal D\otimes_R N))_z\in D^{\le d}(R)\)
for every finite module \(N\), independently of the filtration length.
Thus the error on the geometric side lies in
\(D^{\le d-n-1}(R)\). In every fixed degree both errors disappear for
large \(n\), giving
\begin{equation}\label{geom:coefficient-change}
 P\otimes_R^L k_0\simeq
 \bigl(Ri^!(\mathcal D\otimes_R k_0)\bigr)_z.
\end{equation}
The right-hand side is bounded. A minimal bounded-above free resolution
of \(P\) has differentials in \(\rad R\); tensoring with \(k_0\) kills
them. Boundedness in \eqref{geom:coefficient-change} forces all but
finitely many terms of that resolution to vanish. Hence \(P\) is
perfect. Extension of coefficients preserves the resulting finite
projective complex.
\end{proof}

\begin{lemma}\label{geom:cross-perfect}
For all \(v,x\in W\) and \(\phi\), the complex
\(j_v^!\Delta_{x,\phi}\) is perfect over \(kA_v\).
This remains true when the calculation is performed in an invariant
open containing the two strata.
\end{lemma}

\begin{proof}
Suppress shifts and put \(S=(A_v)_p\), embedded in \(\mathcal H\) as
above; it fixes \(\dot v\). It suffices to restrict to \(S\): since
\([A_v:S]\) is prime to \(p\), a complex of \(kA_v\)-modules is a
retract of the induction of its restriction to \(S\).
Changing equivariance from \(\mathcal H\) to \(S\) is smooth pullback.

Form the finite affine quotient \(\pi:\mathbf G\to Y=\mathbf G/S\).
The maps to the quotient and to the classifying stack give a
representable finite morphism
\[
 \Pi:[\mathbf G/S]\longrightarrow Y\times BS.
\]
Indeed its pullback along the trivial-torsor atlas
\(Y\to Y\times BS\) is \(\pi\). Bruhat strata descend to locally
closed subsets \(Y_x\). Give them reduced structure, which has no
effect on \'etale sheaves. Regard the finite pushforward
\(\mathcal D=\Pi_*(j_{x,!}E_{x,\phi})\) as a sheaf on \(Y\) of modules
over \(R=kS\). This is an ordinary sheaf: extension by zero is exact
and finite pushforward has no higher ordinary cohomology. The
identification with ring-valued sheaves follows from the finite
\'etale nerve \(Y\times S^\bullet\) of the displayed atlas.

Every geometric stalk of \(\mathcal D\) is projective over \(kS\).
Over \(Y_x\), a stalk is induced from the stabilizer in \(S\) of one
point of the corresponding orbit in \(\mathbf G_x\). Such a stabilizer
injects into \(A_x\): torus projection is unchanged by conjugation in
\(\mathcal H\), and its unipotent kernel has no nontrivial finite
\(p\)-subgroup since \(r\ne p\). The stalk coefficient is consequently
the restriction of \(E_{x,\phi}\) to a subgroup of \(A_x\), which is
projective; induction to \(S\) preserves projectivity. Stalks outside
\(Y_x\) are zero.

For \(i:Y_v\hookrightarrow Y\), proper base change and localization
commute \(Ri^!\) with \(\Pi_*\). At \(\pi(\dot v)\) the reduced
fiber is the single fixed point \(\dot v\), so the resulting stalk,
with its \(S\)-action, is exactly the restriction to \(S\) of the
desired cross-costalk. The algebra \(kS\) is commutative local.
Lemma~\ref{geom:flat-costalk} applies after finite descent and proves
perfection. All steps commute with restriction to an invariant open.
\end{proof}

\subsection{Uniformity without a bound on the finite torus}

We next bound the number of indecomposable projectives needed in the
cross-costalks. Their vector-space dimensions are not suitable for
this purpose, because \(\dim E_{x,\phi}=|(A_x)_p|\) can grow.

A rank-one multiplicative Kummer sheaf on a torus means a rank-one
summand of the direct image under an isogeny of degree prime to \(pr\),
with its multiplicative structure. For a fixed \(x\), pullback of the
stratum coefficients to \(\mathbf G_x\) has the following description.
Let
\[
 r_x:\mathbf G_x\longrightarrow\mathbf T,
 \qquad u\dot x b\longmapsto t_b,
\]
where \(u\in\mathbf U\cap{}^{\dot x}\mathbf U^-\) and \(t_b\) is the
torus component of \(b\). There are pairwise distinct multiplicative
Kummer sheaves \(\mathcal L^x_\phi\) on \(\mathbf T\), as \(\phi\)
varies, such that
\begin{equation}\label{geom:torus-coefficients}
 L_{x,\phi}|_{\mathbf G_x}=r_x^*\mathcal L^x_\phi,
 \qquad
 E_{x,\phi}|_{\mathbf G_x}
 =r_x^*(\mathcal L^x_\phi\otimes\mathcal P_x).
\end{equation}
Here \(\mathcal P_x\) is the direct image of the constant sheaf under
a torus isogeny with kernel \((A_x)_p\).
Indeed the torus-coordinate equation in the orbit map is a torus Lang
isogeny with kernel \(A_x\). Factoring it into its \(p\)- and
\(p'\)-parts proves \eqref{geom:torus-coefficients}. A character of the
kernel gives a trivial local system only when it is trivial, because
the covering torus is connected. This proves the asserted distinctness.

We record explicitly the geometric estimate for the rank-one part.

\begin{lemma}\label{geom:kummer-stalk}
Fix the root datum and \(v,x\in W\). Let \(V_v\subset\mathbf G/\mathbf B\)
be the translated opposite big cell through \(\dot v\mathbf B\), with
its natural section \(s:V_v\to\mathbf G\), and put
\(f=r_x\circ s\) on \(C_x\cap V_v\). For every rank-one multiplicative
Kummer sheaf \(\mathcal L\) on \(\mathbf T\) of order prime to \(pr\),
the sum of the dimensions of the stalk cohomology of
\[
 Rj_*(f^*\mathcal L),
 \qquad j:C_x\cap V_v\hookrightarrow V_v,
\]
at any point of \(C_v\) is bounded by a constant depending only on the
root datum. Its restriction cohomology sheaves on \(C_v\) are constant.
\end{lemma}

\begin{proof}
The group \(\mathbf U_v=\mathbf U\cap{}^{\dot v}\mathbf U^-\) acts
simply transitively on \(C_v\). It preserves \(V_v\), the section
satisfies \(s(uz)=us(z)\), and the right torus coordinate is invariant
under left \(\mathbf U\). Thus the entire construction is
\(\mathbf U_v\)-equivariant. Its restriction cohomology on \(C_v\) is
constant, and it suffices to work at \(\dot v\mathbf B\).

Choose a reduced expression \(x=s_1\cdots s_d\). The associated
Bott--Samelson variety \(\widetilde Z_x\) is smooth and proper over
the Schubert closure \(\overline C_x\), is an isomorphism over \(C_x\),
and has a simple normal-crossing boundary with \(d\) distinguished
divisors. In the gallery description these divisors are the conditions
that successive flags are equal. Restrict this proper morphism over
\(V_v\). Its open complement is \(C_x\cap V_v\); write \(\widetilde j\)
for this open immersion and \(\pi\) for the restricted proper morphism.
Then
\[
 Rj_*f^*\mathcal L=R\pi_*R\widetilde j_*f^*\mathcal L.
\]

We describe the cohomology sheaves of the star-extension on a boundary
stratum. A character of a split torus modulo an integer \(n\) prime
to \(pr\) defines its Kummer sheaf by taking an \(n\)-th root of a
torus monomial. The pullback monomial is a unit on the open complement.
In \'etale coordinates near a crossing it has the form
\[
 u\,z_1^{a_1}\cdots z_t^{a_t},
\]
where \(u\) is a unit and the \(z_i\) are parameters for the boundary
divisors. Take \(n\)-th roots of the parameters and of \(u\).
The constant-coefficient cohomology of the punctured crossing is the
exterior algebra on its \(t\) parameter loops. Coordinate deck
translations act trivially on this cohomology. Taking a character
summand is exact because \(p\nmid n\). It follows that the
star-extension vanishes on this stratum if any of its local monodromy
characters is nontrivial. Otherwise its cohomology sheaves are the
extending rank-one Kummer sheaf tensored with exterior powers of
\(k(-1)^t\): trivial local monodromy makes the relevant boundary
powers divisible by \(n\), so the finite Kummer system extends across
these divisors after removing the other boundary components.
The loop lines are constant on the stratum: their residues
are indexed by the globally distinguished divisors. In particular the
sum of their ranks is at most \(2^d\), independently of \(n\).

This bounds the ranks contributed by the boundary crossings.
We now bound the cohomology of these coefficients over the proper
fiber, using a filtration with uniformly boundedly many pieces of
the form \(\mathbb A^a\times\mathbb G_m^b\).
For proper base change, consider the fiber of \(\pi\) at
\(\dot v\mathbf B\). Write \(D_1,\ldots,D_d\) for the equality
divisors. Filter the fiber by the closed conditions that at least
\(m\) equalities hold, in descending order of \(m\). Each layer is a
finite disjoint union of its intersections with the exact strata
\[
 D_I^\circ=
 \Bigl(\bigcap_{i\in I}D_i\Bigr)
       \mathbin{\big\backslash}\Bigl(\bigcup_{j\notin I}D_j\Bigr).
\]
Deleting the steps in \(I\) identifies an exact stratum in the fiber
with the fiber of the uncompactified convolution for the remaining,
possibly nonreduced, simple-reflection word, with every remaining
step unequal. Subdivide further by recording the Bruhat position of
every intermediate flag.
Each nonempty piece is a product
\(\mathbb A^a\times\mathbb G_m^b\), with the number of pieces and
\(a+b\) bounded in terms of \(d\) and \(|W|\). Here is the elementary
gallery verification. Starting with the prescribed final flag, choose
predecessors backwards in the relevant minimal-parabolic projective
line. Its two Bruhat cells have dimensions differing by one. For an
endpoint in the shorter cell, the equal predecessor is a point and
the predecessors in the longer cell form \(\mathbb A^1\). For an
endpoint in the longer cell, the equal predecessor is a point, the
unequal predecessors in that cell form \(\mathbb G_m\), and the
predecessor in the shorter cell is a point. The unipotent section of
each Bruhat cell trivializes these alternatives algebraically as the
endpoint varies. Induction gives the displayed products. Recording
the initial position as the singleton cell \(C_1\) enforces the
initial flag without a further equation. To obtain a filtration on
each exact equality stratum, use the morphism recording intermediate
flags into a product of flag varieties. A linear ordering refining
the product Bruhat order has closed lower ideals. Their inverse
images give a filtration with the fixed-position records as its
successive differences. Combining it with the equality filtration
makes localization applicable, compatibly with the boundary strata.
This is the simple-reflection construction in
\cite[Theorem 1.1, Proposition 5.3, Lemma 6.2,
and Sections 6.2--6.3]{HainesPaving}.

On \(\mathbb A^a\times\mathbb G_m^b\), the Picard group is zero and
every invertible regular function is a constant times a Laurent
monomial. The Kummer exact sequence therefore expresses every
rank-one local system of order dividing \(n\) as a torus Kummer system.
A nontrivial such system has zero torus cohomology: pull back to its
finite prime-to-\(p\) torus cover and use the fact that torus
translations act trivially on cohomology. The constant system has
total compact Betti number \(2^b\), by K\"unneth and duality.
Thus all the extending rank-one systems on the pieces have total
compact Betti number at most \(2^b\).

Apply the cohomology-sheaf spectral sequence on each piece, followed
by the finite localization filtration of the proper fiber. The local
rank bound \(2^d\), the number of pieces, and their dimensions give a
bound depending only on the root datum. Proper base change proves
the claimed stalk bound, uniformly in \(r\), \(n\), and \(\mathcal L\).
More explicitly, there are at most \(3^d\) backward records and each
piece has \(a+b\le d\), so \(12^d\) is a sufficient bound.
\end{proof}

Lemma~\ref{geom:kummer-stalk} bounds the geometric contribution of
one rank-one coefficient.  We now sum over the possible characters
on the receiving stratum.  Character matching and the cohomology
of a covering torus will prevent the degree of the Lang cover
from entering this sum.

\begin{proposition}\label{geom:summed-ext}
For every integer \(J\ge0\), there is a constant \(D_J\), depending
only on the fixed root data and \(J\), such that for all \(v,x,\phi\)
and \(0\le j\le J\),
\begin{equation}\label{geom:ext-bound}
 \sum_{\psi\in\Hom(A_v,k^\times)}
 \dim_k\Ext^j_{\mathcal X}
       (j_{v,!}L_{v,\psi},j_{x,!}E_{x,\phi})\le D_J.
\end{equation}
The same bound holds, after increasing the constant if necessary, in
invariant open unions containing the relevant strata. Bounded shifts
of the two arguments are allowed by increasing \(J\).
\end{proposition}

\begin{proof}
First forget equivariance and compute on \(\mathbf G\). The chart
\(V_v\) contains the whole cell \(C_v\). Its section satisfies
\(r_v(s)=1\) on \(C_v\), and it gives a trivialization
\(\mathbf G|_{V_v}\simeq V_v\times\mathbf B\). Since the first
argument is extended by zero from this open, restriction to it
computes the same Ext. Write \(\rho:V_v\times\mathbf B\to V_v\)
for projection, and \(t_b\) for the torus coordinate of \(b\).
On \((C_x\cap V_v)\times\mathbf B\),
\[
 r_x(s(z)b)=f(z)t_b.
\]
By \eqref{geom:torus-coefficients}, the first coefficient is
\(\mathcal L^v_\psi(t_b)\). Tensor both arguments by its inverse and
use adjunction for \(\rho^*\) and \(R\rho_*\).

The module \(E_{x,\phi}\) has a filtration by copies of
\(L_{x,\phi}\). The corresponding rank-one graded terms, after this
twist, are external products with fiber coefficient
\[
 \mathcal L^x_\phi\otimes(\mathcal L^v_\psi)^{-1}.
\]
Their ordinary direct images are extensions by zero from
\(C_x\cap V_v\) of their fiber cohomology. To justify this assertion
for the nonproper projection, use Verdier duality on the smooth
fiber and compact-support K\"unneth for external products. This proves
both restriction base change and vanishing outside the immersed base
piece. The filtration transfers these two properties to the full
coefficient. Its possibly large length is used only for these
properties, not for a dimension estimate.

A nontrivial multiplicative rank-one torus system has zero ordinary
cohomology, by the torus-cover argument in the preceding proof.
Consequently all Ext groups vanish unless
\begin{equation}\label{geom:character-match}
 \mathcal L^v_\psi\simeq\mathcal L^x_\phi.
\end{equation}
There is at most one such \(\psi\), by the distinctness in
\eqref{geom:torus-coefficients}.

In the matching case the coefficient on the immersed product is
\[
 f^*\mathcal L^x_\phi\otimes\mathcal P_x(f(z)t_b).
\]
The algebraic change of variable \(t'_b=f(z)t_b\) identifies its
ordinary fiber integral with
\[
 f^*\mathcal L^x_\phi\otimes R\Gamma(\mathbf T'_x,k),
\]
where \(\mathbf T'_x\to\mathbf T\) is the torus isogeny defining
\(\mathcal P_x\); the unipotent part of \(\mathbf B\) is acyclic.
More explicitly, write this isogeny as \(h\).  The total cover
\(h(t')=f(z)t_b\) is the product of the immersed base with
\(\mathbf T'_x\), by
\((z,t')\mapsto(z,f(z)^{-1}h(t'),t')\).
Thus the fiber integral is constant as a derived complex, without
choosing a lift of \(f\) through the isogeny.
The covering torus has the same dimension as \(\mathbf T\). Therefore
the total dimension of this constant complex is
\(2^{\rk\mathbf T}\), independent of \(|(A_x)_p|\).
The direct image on all of \(V_v\) is its extension by zero, by the
base-change and vanishing just proved.

It remains to bound
\[
 R\Hom_{V_v}(a_!k,j_!f^*\mathcal L^x_\phi),
 \qquad a:C_v\hookrightarrow V_v.
\]
By adjunction this is
\(R\Gamma(C_v,a^!j_!f^*\mathcal L^x_\phi)\).
Verdier duality expresses its coefficient cohomology in terms of the
restriction to \(C_v\) of
\(Rj_*(f^*\mathcal L^x_\phi)^{\vee}\), with only the fixed smooth
dimension shifts. Lemma~\ref{geom:kummer-stalk} bounds these stalks
and makes their cohomology constant on \(C_v\).
Since \(C_v\) is affine space, ordinary cohomology of a constant
sheaf is concentrated in degree zero. The resulting spectral sequence
bounds the total Ext dimension by a root-datum constant.
This proves the summed estimate on \(\mathbf G\).

The estimate on \(\mathbf G\) is independent of the degree of
the Lang cover.  It remains to restore equivariance, retaining
only the bounded range of cohomological degrees in the statement.
Use the smooth atlas \(\mathbf G\to\mathcal X\).
Cohomological descent for Hom gives a spectral sequence whose term
in simplicial degree \(a\) is an Ext group on
\(\mathbf G\times\mathcal H^a\). Equivariance identifies the pulled
back arguments with projection pullbacks from \(\mathbf G\).
K\"unneth and smooth base change multiply the preceding Ext groups
by \(H^*(\mathcal H^a,k)\). The inputs before shifts are ordinary
sheaves, so these Ext degrees and the simplicial degrees are
nonnegative. Only \(a\le J\) can contribute to total degree at most
\(J\). As a variety, \(\mathcal H\) is a product of a fixed-dimensional
torus and affine space; its relevant Betti numbers, and hence those of
these finitely many powers, depend only on the root datum and \(J\).
Taking dimensions in the spectral sequence proves
\eqref{geom:ext-bound}. For invariant opens the first argument is
extended by zero to the ambient space, so open adjunction and base
change give the same calculation. Negative Ext degrees of ordinary
sheaves are zero, which also proves the assertion about bounded shifts.
For example, writing \(t=\rk\mathbf T\) and
\(d=\max_{w\in W}\ell(w)\), the generous choice
\[
 D_J=2^t12^d\sum_{a=0}^J2^{at}
\]
bounds the sum in every degree at most \(J\).
\end{proof}

\subsection{Gluing with a bounded number of standard factors}

We now combine perfection with the summed Ext estimate.
Perfection reduces each gluing obstruction to two projective terms;
the estimate bounds their multiplicities and hence the length of
the standard flag constructed at the next stratum.

\begin{proposition}\label{geom:bounded-tiltings}
There is a constant \(L\), depending only on the fixed root data,
such that for every \((x,\phi)\) there is an object
\(\mathcal T_{x,\phi}\) with both a standard and a costandard flag.
A specified standard flag has length at most \(L\), contains
\(\Delta_{x,\phi}\) once, and has all remaining factors on strictly
smaller Bruhat strata. Thus the standard multiplicity matrix is
triangular, with identity diagonal blocks.
\end{proposition}

\begin{proof}
Choose a total ordering refining descending Bruhat order. Its initial
unions are invariant opens. When the cell \(x\) is added, start with
the closed pushforward of \(E_{x,\phi}[\ell(x)]\), zero on the
previous cells. This has both flags. Suppose the object \(\mathcal M\)
has been constructed on an old open and \(v\) is the next stratum,
closed in the enlarged open. Write \(j\) for the old open immersion
and \(i\) for the closed stratum immersion, and put
\[
 K=i^!j_!\mathcal M[-\ell(v)].
\]
The standard flag of \(\mathcal M\) and Lemma~\ref{geom:cross-perfect}
make \(K\) perfect over \(kA_v\).
The object \(j_!\mathcal M\) is perverse after atlas pullback and the
fixed \([\dim\mathbf B]\) shift, because it has a standard flag.
Perverse cosupport on the smooth stratum gives \(K\in D^{\ge0}\).
Likewise \(Rj_*\mathcal M\) is perverse because it has a costandard
flag. The localization triangle and \(i^!Rj_*=0\) give
\[
 i^*Rj_*\mathcal M[-\ell(v)]\simeq K[1].
\]
Perverse support puts the left-hand side in \(D^{\le0}\).
Thus \(K\) has cohomology only in degrees \(0,1\).

The algebra \(kA_v\) is split symmetric, hence self-injective. Choose
a minimal bounded complex of projectives representing \(K\).
If its first nonzero term were below degree zero, vanishing of that
cohomology would make its differential injective; self-injectivity
would split this injection, contradicting minimality. If its last
term were above degree one, the preceding differential would be a
surjection and would split by projectivity. Consequently
\[
 K\simeq[P^0\xrightarrow{\delta}P^1],
\]
with the terms in degrees zero and one.
Let \(m^q_\psi\) be the multiplicity of \(E_{v,\psi}\) in \(P^q\).
The terms are also injective, so this complex computes
\(R\Hom_{kA_v}(L_{v,\psi},K)\). Its differential is zero on the
socle of every indecomposable injective summand: a map nonzero on
that essential simple socle would be injective, and hence split,
again contradicting minimality. Since the socle of \(E_{v,\psi}\)
is \(L_{v,\psi}\), it follows that
\begin{equation}\label{geom:multiplicity-ext}
 m^q_\psi=
 \dim_k\Ext^q_{kA_v}(L_{v,\psi},K),\qquad q=0,1.
\end{equation}

Let \(N\) be the old standard-flag length. Adjunction turns the
right-hand side of \eqref{geom:multiplicity-ext} into
\[
 \dim_k\Ext^{q-\ell(v)}
       (j_{v,!}L_{v,\psi},j_!\mathcal M).
\]
For a standard factor indexed by \((y,\eta)\), this is the Ext group
between unshifted ordinary sheaves in degree
\(q+\ell(y)-\ell(v)\).  If
\(d=\max_{w\in W}\ell(w)\), these degrees lie between \(-d\)
and \(d+1\).  Negative Ext degrees of ordinary sheaves vanish,
so we may use \(J=d+1\) in Proposition~\ref{geom:summed-ext}.
Taking long exact sequences along the standard flag gives a
constant \(D\) such that
\begin{equation}\label{geom:glue-size}
 \sum_\psi m^q_\psi\le DN\qquad(q=0,1).
\end{equation}
Computations inside the enlarged open agree with the previous
estimates by open extension adjunction.

We have bounded the two projective terms of the obstruction.
The next extension uses \(P^1\) to add standard factors and
\(P^0\) to supply the matching costandard restriction.
The triangle
\(K\to P^0\to P^1\to K[1]\) gives by adjunction a map
\(i_*P^1[\ell(v)]\to j_!\mathcal M[1]\). Let \(\mathcal M'\) be the
corresponding extension, so that
\[
 j_!\mathcal M\longrightarrow\mathcal M'
 \longrightarrow i_*P^1[\ell(v)]
 \longrightarrow j_!\mathcal M[1].
\]
Its restrictions satisfy
\(j^*\mathcal M'=\mathcal M\) and
\(i^*\mathcal M'=P^1[\ell(v)]\), whereas
\(i^!\mathcal M'=P^0[\ell(v)]\). The displayed triangle gives a
standard flag; the localization triangle
\[
 i_*P^0[\ell(v)]\longrightarrow\mathcal M'
 \longrightarrow Rj_*\mathcal M
\]
gives a costandard flag. By \eqref{geom:glue-size}, the new standard
length is at most \((1+D)N\).

There are at most \(|W|-1\) extension steps, so
\(L=(1+D)^{|W|-1}\) suffices. If the newly added stratum is outside
\(\overline{\mathcal X_x}\), its obstruction complex is zero and
nothing is added. Thus the support stays in
\(\overline{\mathcal X_x}\), and the original
factor on \(x\) occurs exactly once. Taking a maximum over the finite
collection of root data proves the proposition.
\end{proof}

\subsection{Projective representations from the correspondence}

The objects just constructed have bounded flags.  The next step
turns them into projective group representations: the two flags
force their character-functor images into degree zero, while
Rickard's cohomology complex supplies projectivity.  For the
degree bounds we first need affineness in a range uniform over
the fixed data.

The ample-boundary method goes back to Deligne--Lusztig
\cite[Sections~9.5--9.7]{DL}.  For ordinary Frobenius, an
affineness result for all field sizes is given in
\cite[Corollary~3.12]{BrosnanHongLee}; the argument here includes
the specified exceptional diagram isogenies and only needs a
uniform large-parameter range.

\begin{lemma}\label{geom:affineness}
For all sufficiently large \(Q\), uniformly over the specified root
data, maps \(\theta\), and \(x\in W\), the Deligne--Lusztig variety
\[
 X(x)=\{g\mathbf B:g^{-1}F(g)\in\mathbf G_x\}
\]
is smooth affine of dimension \(\ell(x)\).
\end{lemma}

\begin{proof}
Since \(dF=0\), the Lang map is \'etale. Its inverse image of
\(\mathbf G_x\) is a \(\mathbf B\)-bundle over \(X(x)\), which proves
smoothness and the dimension formula.

Here is the ample-line-bundle argument for affineness, including the
exceptional diagram isogenies. On \(\mathcal B=\mathbf G/\mathbf B\)
use the convention
\[
 \mathcal L_\lambda=
 \mathbf G\times^{\mathbf B}\overline{\F}_r{}_{-\lambda};
\]
strictly dominant characters define ample bundles. Choose such an
integral \(\lambda\), which exists for every root datum. On
\(\overline C_x\), the extremal-coordinate functional of weight
\(x\lambda\) in a highest-weight representation gives a
\(\mathbf B\)-semi-invariant section of \(\mathcal L_\lambda\), of
weight \(-x\lambda\), nonzero precisely on \(C_x\). To check its
vanishing on the boundary, along a saturated Bruhat step
\(y<ys_\alpha\) one has
\[
 y\lambda-ys_\alpha\lambda
 =\langle\lambda,\alpha^\vee\rangle y\alpha,
\]
a positive multiple of a positive root. Thus for \(y<x\),
\(y\lambda-x\lambda\) is a nonzero sum of positive roots. Positive
unipotents raise weights, so the coordinate of weight \(x\lambda\)
vanishes on \(C_y\) and is nonzero on \(C_x\). The same calculation
gives the claimed semi-invariance.

The closure of the diagonal relative-position orbit of \(x\) in
\(\mathcal B\times\mathcal B\) is
\(\mathbf G\times^{\mathbf B}\overline C_x\).
Twisting the preceding section by the first-factor line of weight
\(x\lambda\) makes it descend to a section of
\(\mathcal L_{-x\lambda}\boxtimes\mathcal L_\lambda\) whose nonzero
locus is the relative-position orbit itself. Pull back along the graph
\(u\mapsto(u,F(u))\). On the projective inverse image of the orbit
closure its line bundle is the restriction of
\[
 \mathcal L_{F^*\lambda-x\lambda},
 \qquad F^*\lambda=Q\theta^*\lambda.
\]
The pullback \(\theta^*\lambda\) is strictly dominant: the induced
map on flag varieties is finite, and finite pullback preserves
ampleness. This includes inseparable exceptional isogenies. Therefore
\(Q\theta^*\lambda-x\lambda\) is strictly dominant for all
sufficiently large \(Q\), with one threshold for the finitely many
possibilities. The nonzero locus of a section of an ample line bundle
on a projective scheme is affine, by passing to a very ample power.
Here it is \(X(x)\). For a torus the flag variety is a point.
Compare \cite[Sections 9.5--9.7]{DL}.
\end{proof}

Let \(\mathbf B_F^{\mathrm{gr}}\) and
\(\mathbf G_F^{\mathrm{gr}}\) denote the graphs of \(F\), and use
\[
 \mathcal X\xleftarrow{\ \beta\ }
 [\mathbf G/\mathbf B_F^{\mathrm{gr}}]
 \xrightarrow{\ \alpha\ }
 [\mathbf G/\mathbf G_F^{\mathrm{gr}}]\simeq B\Gamma.
\]
The equivalence on the right is Lang's theorem. The map \(\beta\)
is smooth, with unipotent fiber
\(\mathcal H/\mathbf B_F^{\mathrm{gr}}\simeq\mathbf U\), and
\(\alpha\) is proper with flag-variety fiber.
With our ordinary stratum inflation, put
\(\mathsf{Ch}=R\alpha_!\beta^*\), without an additional shift.
Pullback to a point of \(B\Gamma\) gives directly
\begin{equation}\label{geom:character-functor}
 \begin{aligned}
 \mathsf{Ch}(\Delta_{x,\phi})
   &=R\Gamma_c(X(x),\mathcal E_{x,\phi})[\ell(x)],\\
 \mathsf{Ch}(\nabla_{x,\phi})
   &=R\Gamma(X(x),\mathcal E_{x,\phi})[\ell(x)].
 \end{aligned}
\end{equation}
Here \(\mathcal E_{x,\phi}\) is the local system associated to
\(E_{x,\phi}\) on the finite \'etale \(A_x\)-torsor
\[
 Y(\dot x)=\{g\mathbf U:g^{-1}F(g)\in\mathbf U\dot x\mathbf U\}
 \longrightarrow X(x).
\]
For clarity, the flag fiber is stratified by the condition
\(g^{-1}F(g)\in\mathbf G_x\). Write
\(g^{-1}F(g)=c\dot x(c')^{-1}\) with \((c,c')\in\mathcal H\).
The stabilizer-component torsor maps to \(gc\mathbf U\), and
\((c')^{-1}F(c)\in\mathbf U\) gives the defining equation of
\(Y(\dot x)\). This identifies the coefficients and their left
\(\Gamma\)-action. Smooth base change for \(\beta\) handles star
extensions, and properness of \(\alpha\) gives both formulas in
\eqref{geom:character-functor}. Thus their normalization follows from
the fiber calculation itself.

By Lemma~\ref{geom:affineness} and Artin vanishing,
\[
 \mathsf{Ch}(\Delta_{x,\phi})\in D^{\ge0}(k\Gamma),
 \qquad
 \mathsf{Ch}(\nabla_{x,\phi})\in D^{\le0}(k\Gamma).
\]
Indeed a lisse coefficient shifted by \(\ell(x)\) is perverse on the
smooth affine \(X(x)\). Affine Artin vanishing gives the ordinary
cohomology bound, and duality gives the compact-support bound;
the modular coefficient form follows from the same theorem after
finite descent, as in
\cite[Section~4.0, Theorem~4.1.1 and Corollary~4.1.2]{BBD}.
Both flags therefore put
\(\mathsf{Ch}(\mathcal T_{x,\phi})\) in degree zero. Denote this module
by \(P_{x,\phi}\).

\begin{proposition}\label{geom:projective-images}
The modules \(P_{x,\phi}\) are projective. Let
\(\widehat A_x\) be the ordinary linear characters of \(A_x\), with
\(\xi\mapsto\bar\xi\) their modular reductions, and set
\begin{equation}\label{geom:torus-sum}
 \Psi_{x,\phi}=(-1)^{\ell(x)}
       \sum_{\substack{\xi\in\widehat A_x\\\bar\xi=\phi}}
              R_{T_x}^{\mathbf G}(\xi).
\end{equation}
If \(n_{x,\phi;y,\eta}\) are the multiplicities in the specified
standard flag, the ordinary lift character of \(P_{x,\phi}\) is
\begin{equation}\label{geom:projective-character}
 \chi_{P_{x,\phi}}=
       \sum_{y,\eta} n_{x,\phi;y,\eta}\Psi_{y,\eta},
 \qquad \sum_{y,\eta}n_{x,\phi;y,\eta}\le L.
\end{equation}
These projective characters and the \(\Psi_{x,\phi}\) have the same
linear span in characteristic zero, and that span contains the
regular character of \(\Gamma\).
\end{proposition}

\begin{proof}
At a geometric point \(g\mathbf U\) of \(Y(\dot x)\), the stabilizer
for left translation by \(\Gamma\) lies in
\(\Gamma\cap g\mathbf U g^{-1}\), an \(r\)-group, and hence has order
prime to \(p\). Rickard's equivariant cohomology theorem
\cite[Theorems 3.2 and 3.5]{RickardEtale} applies to this
quasi-projective variety with the action of \(\Gamma\times A_x\).
Restricting the resulting complex to \(\Gamma\), Mackey decomposition
gives summands of permutation terms induced from the intersections of point
stabilizers with \(\Gamma\). These intersections are the \(r\)-groups
just described. Thus the bounded complex is termwise projective
over \(\Gamma\), compatibly with integral, residue, and
characteristic-zero coefficients. The commuting right action of
\(A_x\) selects \(\mathcal E_{x,\phi}\) by the idempotent for the
specified character of \((A_x)_{p'}\). This idempotent is integral
over a splitting extension because its denominator is prime to
\(p\), and its direct summand remains perfect over \(k\Gamma\).
Its ordinary Euler character is exactly
\(\Psi_{x,\phi}\): all ordinary characters lifting \(\phi\) are
selected, and the shift in \eqref{geom:character-functor} supplies
\((-1)^{\ell(x)}\).

Thus every standard image is perfect, and a standard flag makes
\(P_{x,\phi}\) perfect. A finite module of finite projective dimension
over the self-injective algebra \(k\Gamma\) is projective. Lifting
projectives over a complete splitting discrete valuation ring
identifies the projective Grothendieck groups before and after
reduction. Taking the Euler classes of the standard flag proves
\eqref{geom:projective-character}; lifting the extension maps is
unnecessary. Proposition~\ref{geom:bounded-tiltings} makes its
multiplicity matrix triangular with identity diagonal blocks, so the
two collections have the same span.

Summing \eqref{geom:torus-sum} over \(\phi\) puts
\(R_{T_x}^{\mathbf G}(\operatorname{Reg}_{A_x})\) in this span for
every \(x\). In a uniform Steinberg formula
\cite[Corollary 7.14, formula (7.14.2), and Section 11]{DL}, replace
each trivial torus character by
\(\operatorname{Reg}_{A_x}/|A_x|\). The Deligne--Lusztig character
formula \cite[Theorem 4.2 and Section 11]{DL} shows that the resulting
class function agrees with \(\St_\Gamma\) on unipotents, since the
values of a torus-induced character there are independent of the
linear torus character. It vanishes on every element with nonidentity
semisimple part, since every corresponding torus evaluation is at a
nonidentity element. The Steinberg character vanishes on nonidentity
unipotents and has nonzero value at the identity. The resulting
function is therefore
\[
 \frac{\St_\Gamma(1)}{|\Gamma|}\operatorname{Reg}_\Gamma.
\]
This proves the final span assertion.
\end{proof}

\subsection{Projection to a block and the Cartan estimate}

The projective characters now span the regular character, so they
detect every projective indecomposable summand.  Their full norms
can still grow with the tori.  The remaining estimate therefore
uses only the ordinary coordinates in the chosen bounded-defect
block, including after a central quotient.

For a finite group, write \(\langle\ ,\ \rangle\) for the ordinary
character inner product and \(\|\chi\|^2=\langle\chi,\chi\rangle\).
Let \(K(M)\) be a Brauer--Feit bound for the number of ordinary
irreducible characters in any block with defect groups of order at
most \(M\).

\begin{lemma}\label{geom:projected-norm}
Let \(\mathcal I\) be any set of at most \(K\) ordinary irreducible
characters of \(\Gamma\), and let \(\operatorname{pr}_{\mathcal I}\)
be their orthogonal coordinate projection. Then
\[
 \|\operatorname{pr}_{\mathcal I}\Psi_{x,\phi}\|^2\le K|W|^3,
 \qquad
 \|\operatorname{pr}_{\mathcal I}\chi_{P_{x,\phi}}\|^2
       \le KL^2|W|^3.
\]
\end{lemma}

\begin{proof}
Deligne--Lusztig orthogonality gives
\(\|R_{T_x}^{\mathbf G}(\xi)\|^2\le |W|\)
for every linear torus character \(\xi\); see
\cite[Theorem 6.8 and Section 11]{DL}. Hence every individual
irreducible coefficient has absolute value at most \(\sqrt{|W|}\).
Fix an irreducible character \(\rho\). By geometric-series
disjointness and torus duality, the parameters \(\xi\) from a fixed
torus whose Deligne--Lusztig characters contain \(\rho\) have dual
semisimple elements in one geometric conjugacy class. The intersection
of that class with the fixed dual maximal torus is one Weyl orbit.
It has at most \(|W|\) elements. Thus at most \(|W|\) summands in
\eqref{geom:torus-sum} can contribute to the \(\rho\)-coordinate,
irrespective of how many ordinary characters lift \(\phi\).
Consequently
\[
 |\langle\Psi_{x,\phi},\rho\rangle|\le |W|^{3/2}.
\]
Summing squares over \(\mathcal I\) proves the first inequality.
Equation~\eqref{geom:projective-character} and its bound \(L\) on
the number of factors prove the second.
\end{proof}

\begin{proof}[Proof of Theorem~\ref{thm:geometric-cartan}]
First assume \(Q\) exceeds the uniform threshold used above.
Take central coinvariants
\[
 \overline P_{x,\phi}
   =k[\Gamma/Z]\otimes_{k\Gamma}P_{x,\phi}.
\]
These are projective \(k[\Gamma/Z]\)-modules, since tensoring an
exhibited summand of a free module gives a summand of a free quotient
module. The same construction on integral projective lifts commutes
with reduction. In characteristic zero its character is the
coordinate projection onto irreducibles trivial on \(Z\).
This observation applies also when \(p\mid |Z|\).
Project further to \(b\), obtaining actual projective modules
\(Q_{x,\phi}=b\overline P_{x,\phi}\).

The regular-character span in Proposition~\ref{geom:projective-images}
projects to the regular character of \(b\): the coordinates of
\(\operatorname{Reg}_\Gamma\) on characters inflated from
\(\Gamma/Z\) are exactly their degrees, as in
\(\operatorname{Reg}_{\Gamma/Z}\). Characters of projective
indecomposable modules are linearly independent, and every such
module in \(b\) occurs with positive multiplicity in its regular
module. Hence every projective indecomposable module of \(b\) occurs
as a summand of some \(Q_{x,\phi}\). Otherwise its coordinate would
vanish in the span of all their projective characters, contradicting
the regular-character assertion.

Inflate the set \(\Irr(b)\) to \(\Gamma\). It has at most \(K(M)\)
members, so Lemma~\ref{geom:projected-norm} gives
\[
 \|\chi_{Q_{x,\phi}}\|^2\le K(M)L^2|W|^3.
\]
If \(\Phi_i\) is an indecomposable projective character occurring in
\(\chi_{Q_{x,\phi}}\), nonnegativity of ordinary multiplicities gives
\(\|\Phi_i\|\le\|\chi_{Q_{x,\phi}}\|\).
The Cartan entries are
\(c_{ij}(b)=\langle\Phi_i,\Phi_j\rangle\), so Cauchy--Schwarz yields
\[
 c_{ij}(b)\le K(M)L^2|W|^3.
\]
The root data are fixed throughout this estimate; both \(|W|\) and
\(L\) depend only on their specified finite collection. Finally,
the omitted bounded range of \(Q\) contains only finitely many finite
groups and central quotients. Increase the constant by the maximum
of their relevant Cartan entries. This proves the theorem.
\end{proof}
